\chapter{Programs — Whole Program}
%%% generated by the blueprint pipeline from TCSlib.Programs.Light.Compiler.WholeProgram
\section{Overview}

Declarations of the Lean module \texttt{TCSlib.Programs.Light.Compiler.WholeProgram}.

\section{Declarations}

\begin{definition}[Decision verdict]
\lean{Light.verdictOf}
\label{Light.verdictOf}
\leanok
For a Boolean flag $d$ and an integer $r$, define the verdict to be true exactly when
either $d$ is false or $r$ is positive.
\end{definition}

\begin{definition}[WordRAM verdict condition]
\lean{Light.Compiler.Verdict}
\label{Light.Compiler.Verdict}
\leanok
\uses{Light.Compiler.cD, Light.Compiler.cRESULT, Light.verdictOf, WordRAM.AgreeOutside, WordRAM.Cfg, WordRAM.Instr, WordRAM.step, WordRAM.wd}
For a word width $W$, program $code$, Boolean value $dec$, integer $r$, initial memory
$m$, and configuration $cfg$, the execution of $code$ from $cfg$ must terminate with the
verdict determined by $dec$ and $r$. The configuration’s result register must contain
the $W$-bit encoding of $r$, and its memory must agree with $m$ outside the result and
decision registers.
\end{definition}

\begin{theorem}[Execution of the verdict tail]
\lean{Light.Compiler.steps_tail}
\label{Light.Compiler.steps_tail}
\leanok
\uses{Light.Compiler.Verdict, Light.Compiler.cD, Light.Compiler.cRESULT, Light.Compiler.cStack, Light.Compiler.cZERO, Light.Compiler.mainFrame, Light.Compiler.tailCode, Light.verdictOf, WordRAM.AgreeOutside, WordRAM.AgreeOutside.refl, WordRAM.AgreeOutside.update, WordRAM.Cfg, WordRAM.CodeAt, WordRAM.CodeAt.head, WordRAM.CodeAt.tail, WordRAM.InRange, WordRAM.Instr, WordRAM.Steps, WordRAM.Steps.trans, WordRAM.step_accept, WordRAM.step_reject, WordRAM.steps_add, WordRAM.steps_bltz, WordRAM.steps_sub, WordRAM.wd}
\difficulty{5}
Let $W$ be a word width, let $P$ be a program, and let $p$ be a program position such
that the tail code for a Boolean decision $d$ at position $p$ occurs in $P$ starting at
$p$. Suppose a memory $m$ stores the word representing an integer $r$ at the stack
location of the main frame and the word representing $0$ at the zero location, and
suppose $-2^W \le -2r < 2^W$. Then execution from position $p$ with memory $m$ reaches,
in at most three steps, a configuration whose execution terminates with the verdict that
is true exactly when $d$ is false or $r$ is positive; its result location stores the
word representing $r$, and its memory agrees with $m$ outside the result and decision
locations.
\end{theorem}

\begin{definition}[Program start cost]
\lean{Light.Compiler.startCost}
\label{Light.Compiler.startCost}
\leanok
\uses{Light.Compiler.dispPos, Light.Compiler.frameSize, Light.Program}
For a program $P$ and a Boolean flag $d$, define its start cost as the program's
dispatch position for $d$, plus its frame size, plus $13$.
\end{definition}

\begin{definition}[RAM step count]
\lean{Light.ramSteps}
\label{Light.ramSteps}
\leanok
\uses{Light.Compiler.dispPos, Light.Compiler.startCost, Light.Compiler.stepsPerStep, Light.Program}
For a program $P$, a Boolean flag $d$, and a natural number $c$, define the RAM step
count to be the start cost for $P$ and $d$, plus $c$ times the number of steps per step
at the program's dispatch position for $d$.
\end{definition}

\begin{definition}[Stack cell count]
\lean{Light.Compiler.stackCells}
\label{Light.Compiler.stackCells}
\leanok
\uses{Light.Compiler.frameSize, Light.Compiler.mainFrame, Light.Limits, Light.Program}
For a program $P$ and limits with depth $h$, define the stack cell count as the main
frame size plus $(	ext{frame size of }P+1)(h+1)$.
\end{definition}

\begin{definition}[Low-cell bound]
\lean{Light.Compiler.lowCell}
\label{Light.Compiler.lowCell}
\leanok
\uses{Light.Compiler.dispPos, Light.Compiler.stackCells, Light.Limits, Light.Program}
For a program $P$, a Boolean flag $d$, and resource limits, define the low-cell bound to
be the larger of twice the stack-cell count for $P$ under those limits and $25$ plus
four times the dispatcher position for $P$ and $d$.
\end{definition}

\begin{theorem}[Frame allocation bound]
\lean{Light.Compiler.frame_le_stackCells}
\label{Light.Compiler.frame_le_stackCells}
\leanok
\uses{Light.Compiler.frameSize, Light.Compiler.mainFrame, Light.Compiler.stackCells, Light.Limits, Light.Program}
\difficulty{3}
For every program $P$ and limits record $	ext{lim}$, the main frame size plus the frame
size of $P$ plus one is at most the number of stack cells allocated for $P$ under
$	ext{lim}$: $4 + 	ext{frameSize}(P) + 1 \le 4 +
(\text{frameSize}(P)+1)(\text{lim.depth}+1)$.
\end{theorem}

\begin{definition}[Compiler setting construction]
\lean{Light.Compiler.setting}
\label{Light.Compiler.setting}
\leanok
\uses{Light.Compiler.Fits, Light.Compiler.Setting, Light.Compiler.codeAt_bodies, Light.Compiler.codeAt_bodiesCode, Light.Compiler.codeAt_dispatcher_dispPos, Light.Compiler.codeLayout, Light.Compiler.dispPos, Light.Compiler.stackCells, Light.Compiler.two_mul_frameSize_add_le_dispPos, Light.Compiler.width_le_frameSize, Light.Program, Light.compileProgram}
Let $W$ be a word size, let the resource limits, program $P$, designated procedure
index, and Boolean flag be fixed, and suppose the word size fits the limits, dispatcher
position, and stack-cell count required by the compiler. The compiler setting for these
data consists of the program layout and compiled code, with each procedure body followed
by its return sequence before the dispatcher, the dispatcher at its specified position,
and every procedure body of width at most the frame size.
\end{definition}

\begin{theorem}[Initial main statement setup]
\lean{Light.Compiler.start_mainStmt}
\label{Light.Compiler.start_mainStmt}
\leanok
\uses{Light.Compiler.Rel, Light.Compiler.Start, Light.Compiler.codeAt_mainStmt, Light.Compiler.codeLayout, Light.Compiler.frame_le_stackCells, Light.Compiler.mainFrame, Light.Compiler.mainPos, Light.Compiler.setting, Light.Compiler.size_mainStmt, Light.Compiler.stackCells, Light.Compiler.three_le_frameSize, Light.Compiler.two_mul_frameSize_add_le_dispPos, Light.Compiler.width_mainStmt, Light.State, Light.State.Bounded, Light.mainStmt}
\difficulty{4}
Let $W$ be a word width, let $lim$ be a resource limit, let $P$ be a program, let $p_0$
be a procedure index, and let $dec$ be a Boolean. Suppose the word width fits the
limits, the dispatcher position determined by $P$ and $dec$, and the stack capacity
determined by $P$ and $lim$. For every state $σ$ bounded by $lim$ and every word memory
$m$ related to $σ$ at stack position $4$, the statement that calls procedure $p_0$ with
the first two variables as arguments satisfies the starting conditions at depth $0$,
code position $4$, and stack position $4$.
\end{theorem}

\begin{theorem}[Addresses outside the reserved regions]
\lean{Light.Compiler.not_mem_of_far}
\label{Light.Compiler.not_mem_of_far}
\leanok
\uses{Light.Compiler.Constants, Light.Compiler.MainFrame, Light.Compiler.Writable, Light.Compiler.cD, Light.Compiler.cFP, Light.Compiler.cRESULT, Light.Compiler.codeLayout, Light.Compiler.dispPos, Light.Compiler.frameSize, Light.Compiler.frame_le_stackCells, Light.Compiler.lowCell, Light.Compiler.mainFrame, Light.Compiler.setting, Light.Compiler.stackCells, Light.Compiler.two_mul_frameSize_add_le_dispPos}
\difficulty{4}
Let $W$ be a natural number, let $	ext{lim}$ be a record of word, space, and depth
limits, let $P$ be a light program, let $p_0$ be a natural number, and let $	ext{dec}$
be a Boolean. Assume the word size fits the limits, the pool bound
$	ext{dispPos}(P,	ext{dec})$, and the stack size
$	ext{mainFrame}+(	ext{frameSize}(P)+1)(	ext{lim.depth}+1)$. For any integer $a$, if
$a<-	ext{lowCell}(P,	ext{dec},	ext{lim})$ or $	ext{lim.space}\le a$, then $a$ belongs to
none of the following: the union of the constants for pool bound
$	ext{dispPos}(P,	ext{dec})$, the frame pointer and main-frame addresses for frame size
$	ext{frameSize}(P)$, the writable addresses of the generated size configuration at
stack index $	ext{mainFrame}$, and the result and difference-register addresses.
\end{theorem}

\begin{definition}[Bounded memory representation]
\lean{Light.Compiler.Holds}
\label{Light.Compiler.Holds}
\leanok
\uses{Light.Limits, WordRAM.wd}
Let $W$ be a natural-number word width, let $L$ consist of an integer word bound and
natural-number space and depth bounds, let $\mu$ assign an integer to each
natural-numbered address, and let $m$ assign a $W$-bit word to each integer address. The
representation is valid if, for every natural-numbered address $a$, $m(a)$ is the
$W$-bit encoding of $\mu(a)$, and $|\mu(a)|$ is at most the integer word bound in $L$.
\end{definition}

\begin{definition}[Input memory condition]
\lean{Light.Compiler.Input}
\label{Light.Compiler.Input}
\leanok
\uses{Light.Compiler.Holds, Light.Compiler.cARG1, Light.Compiler.cARG2, Light.Limits, WordRAM.wd}
Let $W$ be a word size, let $\mathrm{lim}$ be a triple of resource limits, let
$\mu\colon\mathbb{N}\to\mathbb{Z}$ assign an integer to each natural-numbered address,
and let $m\colon\mathbb{Z}\to\{0,1\}^W$ be a memory of $W$-bit words. For integers $i$
and $j$, the input condition requires that $m(a)=\mathrm{wd}_W(\mu(a))$ and
$|\mu(a)|\leq\mathrm{lim.word}$ for every natural number $a$, that the memory cells at
the first and second argument addresses encode $i$ and $j$, respectively, and that
$|i|,|j|\leq\mathrm{lim.word}$.
\end{definition}

\begin{definition}[Program execution outcome]
\lean{Light.Compiler.Outcome}
\label{Light.Compiler.Outcome}
\leanok
\uses{Light.Compiler.Holds, Light.Compiler.cARG1, Light.Compiler.cARG2, Light.Compiler.cRESULT, Light.Compiler.lowCell, Light.Limits, Light.Program, Light.State}
Let $W$ be a word size, let $\text{lim}$ be resource limits, let $P$ be a program, let
$d$ be a Boolean flag, and let $m,m'$ be memories of $W$-bit words indexed by integers.
For initial and final program states $\sigma,\sigma'$, an outcome means that $m'$
represents the final memory of $\sigma'$ with all values within the word bound; the
reserved result cell encodes $\sigma'.loc(0)$; both reserved argument cells are
unchanged from $m$ to $m'$; and every cell strictly below minus the low-cell bound or at
or above the space limit is unchanged.
\end{definition}

\begin{theorem}[Correctness of program compilation]
\lean{Light.Compiler.compileProgram_correct}
\label{Light.Compiler.compileProgram_correct}
\leanok
\uses{Light.Compiler.Fits, Light.Compiler.Fits.inRange, Light.Compiler.Input, Light.Compiler.Outcome, Light.Compiler.Sizes.inRange_neg_one, Light.Compiler.agreeOutside_startMem, Light.Compiler.codeAt_tailCode, Light.Compiler.codeLayout, Light.Compiler.dispPos, Light.Compiler.frameSize, Light.Compiler.mainFrame, Light.Compiler.mainPos, Light.Compiler.not_mem_of_far, Light.Compiler.rel_startMem, Light.Compiler.setting, Light.Compiler.sim, Light.Compiler.size_mainStmt, Light.Compiler.stackCells, Light.Compiler.startCost, Light.Compiler.start_mainStmt, Light.Compiler.stepsPerStep, Light.Compiler.steps_start, Light.Compiler.steps_tail, Light.Compiler.tailPos, Light.Compiler.three_le_frameSize, Light.Compiler.two_mul_frameSize_add_le_dispPos, Light.Exec, Light.Exec.bounded, Light.Exec.mem_outside, Light.Limits, Light.Program, Light.State, Light.State.Bounded, Light.bounded_frame, Light.compileProgram, Light.frame, Light.mainStmt, Light.ramSteps, Light.verdictOf, WordRAM.AgreeOutside.cell, WordRAM.AgreeOutside.read, WordRAM.AgreeOutside.trans_union, WordRAM.Steps.cast_pos, WordRAM.Steps.trans, WordRAM.exec, WordRAM.exec_mono, WordRAM.exec_of_steps, WordRAM.toInt_wd}
\difficulty{6}
Let $P$ be a program, let $	ext{lim}$ specify word, space, and depth limits, let $p_0$
be a procedure index, and let $	ext{dec}$ be a Boolean. Suppose that the main call to
procedure $p_0$, executed in $P$ at depth $0$ from the state whose local values are
$(0,i,j)$ and whose memory is $\mu$, reaches a state $\sigma'$ with cost $c$. Suppose
the word width $W$ satisfies $2(2\,\text{lim.word}+2)<2^W$, $2\,\text{lim.space}<2^W$,
$4Q<2^W$ for $Q=\mathrm{stackCells}(P,\text{lim})$, and $2N<2^W$ for
$N=\mathrm{dispPos}(P,\text{dec})$. Suppose also that the initial word-RAM memory
represents $\mu$ in every nonnegative cell, every value of $\mu$ has absolute value at
most $\text{lim.word}$, and its two argument cells contain $i$ and $j$, respectively,
with $|i|,|j|\le\text{lim.word}$. Then running the compiled program from position $0$
for $\mathrm{ramSteps}(P,\text{dec},c)$ steps returns true exactly when $\text{dec}$ is
false or $0<\sigma'.\mathrm{loc}(0)$. Moreover, some resulting word-RAM memory
represents the final memory of $\sigma'$ in every nonnegative cell, has
$\sigma'.\mathrm{loc}(0)$ in its result cell, preserves both argument cells, and agrees
with the initial memory at every address below
$-\mathrm{lowCell}(P,\text{dec},\text{lim})$ or at least $\text{lim.space$.}
\end{theorem}
